\section{Core Technologies}
\label{sec-core}

In the following paragraphs, I enumerate the technologies
we plan to work on to tackle the challenges, 
beginning with our efforts in synchronizing media streams.

\subsection{Temporal Indexing for Scene Synchronization}

Content interactivity allows a viewer to change the content
of a program locally. For instance, the viewer can change
the attributes (e.g., colors, positions and speed)
of the objects in a scene. 
Content interactivity can be enriched if the TV is connected
to a bi-directional Internet back-channel.
Using distance education as an example,
a student can interact with a 3D globe
during a world geography lecture.
The instructor can quiz the students on-line to understand 
how the students absorb the lecture materials and can make real-time
adjustments.

To support content interactivity, the receiver must synchronize
the rendering of the data objects in a scene. 
%Fortunately, much low-hanging fruit 
%in this area is ready to be plunked. 
This knowledge is already available.
For example, MPEG-4 \cite{Mp4} has proposed standards to
separate interactive objects from the scenes, to change
visual or acoustic properties of an object, and to
specify the spatio-temporal relations between objects.
Many studies have also proposed synchronization models
and languages \cite{Hyp92,Co96,Gr95,Ha96,HFK95,Hkm97}.
However, neither MPEG-4 nor these studies 
address {\em data organization}
and {\em QoS adaptivity} issues, issues I intend to investigate.

\begin{figure}[t]
\epsfysize=140pt
\centerline{\epsffile{fig-scene-desc.eps}}
\caption{An Example Scene Description.}
\label{fig-scene-desc}
\end{figure}
                  
Figure~\ref{fig-scene-desc} shows a schematic example of these issues,
an MPEG-4 scene description of a simplified geography
lecture. The scene at its top level is divided into four
objects: a 2D background, furniture, a white board and 
an instructor.
The furniture ``object'' consists of a desk, a chair and a globe
``subobject.''  The instructor object consists of a 
voice and a sprite subobject.  
Of these objects, a student 
can interact with the white board and the globe.
What this means is that the 2D background and the instructor 
are broadcast via satellite, and the furniture and white board
are delivered via the Internet. The receiver faces the
challenge of synchronizing
these two streams in both spatial and temporal dimensions.

To simplify our discussion,
let us focus on the temporal synchronization.
The receiver must properly synchronize two streams in the temporal
dimension.  If streams are delivered and
rendered in real-time, the temporal synchronization may not be
a major problem since the receiver can render the objects in
the streams as they arrive. However, if the viewer issues a 
pause, the streams may be stored on the local disk, and how 
they are stored on disk will affect their subsequent
retrieval. 
%For example, the 
%temporarily related objects can be stored contiguously
%on disk. This approach, however, may not be feasible
%if objects appear in many scenes.
%Various data organization policies will be experimented with
%to discover rules for storing temporally related data.

When a data object
is referenced, it and all its temporally related objects should be
found quickly, and, once requested, be retrieved efficiently.
Indexing temporally related objects can be tricky.
In principle, all temporally related object should be
stored contiguously on disk to minimize the disk latency
necessary to retrieve them.
However, since an object may be shared in many scenes,
replication may not always be desirable. One alternative is
storing shared objects in the middle tracks so their
access latency is constrainted. Another alternative
is storing shared objects on the outer tracks to take
advantage of their higher data transfer rates.
We will experiment with various data placement policies to
examine if some rules of thumb can work most of the time.

It seems that building temporal indexes at the scene level is
appropriate. Once a scene can be located, the
scene description can tell where its data objects
reside. I believe that the combination of hashing and XML
may be an appropriate structure for indexing scenes \cite{PBC99}.
A hash table provides fast scene lookup; XML
provides the flexibility to describe semistructured
objects in a scene. We will investigate how to
deal with efficient updates (primarily necessary because of changes
of the object's location)
and queries on XML structures \cite{Abi97,Bun97,MW99}.

One useful feature that the receiver can support is video
editing (changing the temporal sequence of 
scenes). A video editing tool can combine scenes from
different videos into a customized clip. Supporting
video editing introduces additional scene management issues.
When a scene is shared by two or more clips, the scene
should not be copied entirely. A scene is copied only
when its substantial content is modified. Otherwise,
we may want to create a shell scene, which maintains
links to its parent scene and keeps only those objects that it
modified.
              
\subsection{Caching and Prefetching Techniques}

Viewing interactivity allows viewers to select their own
temporal view of the TV content in real time.
A temporal view is created by a time-shift VCR-like 
operation---pause, replay or fast-forward. Extending a DRAM cache
by adding a disk increases the cache size and thereby
enables the support of these time-shift operations economically.
However, adding a disk to a DRAM cache also increases the complexity of
the receiver's buffer management.
This increased complexity is primarily introduced by the fact that
video and audio data must be made available continuously
to meet {\em real-time} playback deadlines \cite{ex9}, or playback 
interruptions occur. In addition, to achieve acceptable 
interactive performance,
the buffer manager must make the requested data available
in main memory with {\em minimum initial latency}, e.g.,
when a viewer issues a replay, the response time should be very short.
Satisfying these two potentially conflicting performance 
objectives with minimum DRAM (and hence at minimum cost) 
is a challenging design task.
For instance, in the geography lecture example, a
student my click on the globe object on the desk to request a
video clip about the earth, while the instructor
is speaking. We do not want to interrupt the
continuation of the audio track but we also want
to render the video clip as soon as possible.
To deal with these conflicting goals will be a challenge.

Much work has been done at the server side to schedule IOs
for retrieving
multimedia data \cite{ex121,ex86,ex8-1,ex18,ex302,ex10}.
It is important to note that the {\em receiver-side}
analysis differs from these more
``traditional'' server-side analyses.
In particular, in a server environment, the server has full control of
the data generation rate, while at the receiver, data arrival (from
the network) is unpredictable.
In other words, the server schedules events while the receiver
reacts to events.

To achieve minimum initial latency,
the buffer manager
must be able to service an interactive request
promptly. Since the disk IO is the bottleneck of
the performance, the objective is to schedule IOs
in the right order so that the disk arm on the one hand can
quickly service the newly requested
interactive task and on the other hand can fulfill
the real-time constraints of other data objects.
                                         
Three techniques for making IO scheduling more flexible can
be explored. They are: {\em data
prefetching}, {\em caching}, and {\em data placement}.
When the disk bandwidth is under-utilized, the disk
can be scheduled to perform prefetching
for continuous data so that when
an interactive request arrives, a sufficient amount of
continuous data has
been staged in main memory to meet their real-time
requirements. This gives the disk arm
freedom to service the interactive request immediately.
To minimize the DRAM requirement (and hence reduce the
cost of the receiver), the buffer manager
should minimize disk latency, which requires implementing
intelligent disk-scheduling and data-placement policies.
We will evaluate the tradeoffs between
various factors (e.g., cost, speed and complexity) and devise
efficient and cost-effective buffer management strategies
for the receiver to meet the low initial latency goal 
under real-time data constraints.
                              
\subsection{QoS-Adaptive Buffer Management}

Adding complexity to the problems, the buffer manager must be
{\em QoS adaptive} and must be able to scale gracefully in the following
two areas:
\begin{itemize}
\item Resources:
As mentioned earlier, the
available hardware resources can differ from one receiver
to another. The local resources may not be adequate
to support a particular interactive scenario.
For instance, interacting with an object in a scene 
may require downloading a high
resolution video clip, which the processor cannot decode
in real-time.  Therefore, it is critical for the resource 
manager to be adaptive to
the resource constraints and to degrade, if degradation is
unavoidable, in a graceful manner.  
\item QoS requirements:
The QoS preference can be different from one viewer to another.
For instance, some students watching a televised lecture 
may prefer seeing the write board
to the fidelity of the video, or vice versa.
\end{itemize}

The degradation of content can happen on two different
levels: on the object level and on the scene level.
                                                 
\begin{itemize}
\item On the object level, a layered coding scheme can
be used so that only selected layers are decoded to achieve
the targeted QoS.  Much effort has been devoted to developing a
layered coding scheme that supports scalable QoS \cite{ex220,ex62},
and there have also been studies on QoS adaptivity
in the network transport layer \cite{ex56,ex96,ex95}.
We will evaluate these technologies and choose the ones
that suit our needs.

\item On the scene level, some data objects
that are occluded can be kept from
being processed. The remaining objects can be scheduled
based on a combination of the viewer's QoS specification
and the delay functions of the data objects. The goal
of the scheduler is to maximize the total QoS
under the resource constraints.
Suppose that in the geography lecture example the
system resources cannot render both the white board object
and the instructor object. In this case the resource manager
renders the object that receives a higher QoS assignment   
from the student.
\end{itemize}

The challenge posed by all the above issues---caching, data prefetching
and placement, QoS-adaptivity, initial latency, real-time
data constraints and cost---will be to
devise a resource management policy that satisfies all
requirements at the same time.
To achieve this aim, we plan to
first develop effective techniques for caching, data
prefetching and placement, and scheduling, and then to
combine individual techniques to achieve our global
objective: maximizing QoS. We plan to
formulate the resource management plans quantitatively.
Our preliminary idea is to develop a model that considers all
constraints together with QoS requirements
and that can be solved efficiently in real time
for a schedule that achieves maximum total QoS.
                                                     
\subsection{Media Query Processing}

To support effective time-shift operations, the receiver needs
a {\em search engine} that supports fast lookup and retrieval
of the viewer's desired video clip. For instance, a
viewer may want to replay a particular section of a
lecture in which the history of television is discussed.
The search engine looks up the indexes to find
scene descriptors that best match the key words ``television
history'' and returns
a list of scenes in a ranking order.
Another way to specify a query is content-based. The viewer can draw
an object and its trajectory \cite{Ccmsz97}, and 
the search engine finds
and returns the scenes that have a similar object 
moving in the same trajectory. 

For each video, the search engine builds its scene descriptors and stores
these descriptors in the database.  This process can be divided into
two major steps: {\em feature selection} and {\em indexing}. 

{\it Feature selection}
has enjoyed much attention lately, as is witnessed by the
large number of recent articles on it.
For example, \cite{Ccmsz97,Hjw95,Mmz95,Mc96,Zks93} propose
feature extraction techniques for video, \cite{Bc94,Pfe96} for audio
and \cite{img11,img12,img13,Pps96,Pzm96,Hkm97} for images.
In addition to evaluating and using some of these
existing approaches, we will experiment with 
a {\em variance control model}
that decides how fine the features used should be 
based on 1) the user's definition of a ``match'' 
and 2) the size of the database.  Take image searching as an example.
If one wants to find near-copies of an image,
the search engine should extract fine details
from the query image and perform a close match.
If one wants to find  
images of similar colors that are not necessarily
at the same spots, the seach engine
can extract a coarser color distribution from the query image
to perform the lookup.  Consider the
size of the database. While a small database 
(the typical database of a receiver) requires a 
small set of features to tell media objects apart, a
large image database (the typical database of a content provider)
requires a large set of features.
Intuition tells us that the selection of features
during the indexing and query processing steps
should depend on the semantics of
the search, the size of the database, and the 
available resources.  We will investigate
if this variance-controlled approach can be fruitful.

Content-based retrieval
is inherently expensive, especially when
it needs a large number of dimensions to characterize the
target object (e.g., an image or a video scene).
There is simply no way to
build an index on disk such that {\em all} nearest neighbors to {\em
any} query point are physically near on disk.
Fortunately, in many cases it is
sufficient to perform an {\em approximate search} that returns many
but not all nearest neighbors~\cite{AM94,Cla94,IM98,Jk97,KOR98}.  (A
feature vector is often an approximate characterization of an object,
so we are already dealing with approximations anyway.)
I have recently proposed a
clustering/indexing combined scheme that achieves approximate
similarity search with high efficiency \cite{Lch99}.
This proposed clustering/indexing combined scheme should
provide a solid foundation for continued work,
but many challenges remain. We are, for one, not
convinced that the Euclidean distance
is the only way to measure similarity for multimedia data.
Thus, we are investigating other similarity measures
(e.g., Hamming distance, etc.). For another,
we think similarity should be personalizable.
We will thus investigate new measures for similarity.

Finally, a challenging issue about building indexes for video scenes
is that some scenes may have missing features and that different
scenes may be best represented by different sets of
features.  For instance, the video frames in a lecture and in a news program
may provide little useful information about the content
of the program if all they show is a ``talking head.''
In this case, using their audio tracks
and text captions may better tell the content. For scenes that
have motion, on the other hand, audio tracks and captions may not be useful.
Rather, the motion vectors in the inter-coded frames
can tell the speed and trajectory
of the moving object. We plan to investigate if
assigning different weights to different features may 
help promote those features with greater significance.
We will also investigate techniques that can deal
with missing features when indexes are built.


